\section{The overlay/solid contract}\label{sec:overlay}

This chapter explains the deepest design idea in the toolkit -- why a
solutions copy can show every answer without moving a single page break --
and the one switch that controls it.

\subsection{Reserved space: \cs{workspace}}

Assessments and lectures need copious blank room for handwriting, which is
not the LaTeX norm. \cs{workspace} reserves it \emph{explicitly and visibly
in the source}, as a leaf command between content (never a wrapper around
it), so the source stays readable and the space can be tuned in the editor:

\begin{manexsource}
\workspace            % 8 lines (the default)
\workspace[12]        % 12 lines
\workspace[halfpage]  % about half the text block
\workspace[fullpage]  % the whole text block
\workspace*           % fill to the bottom of this page
\end{manexsource}

The unit is a \emph{line of handwriting} (\cs{HandwritingLine}, about
0.9\,cm) because ``about ten lines'' is how you actually think about writing
room -- centimetres are not. Retune \cs{HandwritingLine} once and
every line-count workspace follows.

One deliberate simplicity: a workspace is plain vertical space, with
\emph{no carry-over}. If it exceeds the room left on the page, the overflow
is dropped and the next content starts the next page -- the gap does not
reappear overleaf. The rare page where that matters is re-spaced by hand,
rather than letting a half-page gap silently wander onto the next page. (A
consequence: like any top-of-page glue, a workspace at the very top of a
page is discarded -- put content above it, or use \cs{workspace*}.)

\cs{FreshPage} completes the spacing toolkit with collapsible page breaks:
bare \cs{FreshPage} breaks to a fresh page that content then flows onto;
\cs{FreshPage}\texttt{[blank]} ships a dedicated empty page;
\texttt{[watermark]} stamps that page with a faint diagonal ``Blank page for
extra work.'' Plain \cs{newpage} remains untouched, for the rare break that
must never collapse.

\subsection{Overlay mode: the default}

Here is the contract. In the default \emph{overlay} mode, the gated boxes --
\env{solution} and \env{InstructorBox} -- render at \textbf{zero height},
painting \emph{downward into the workspace reserved below them} instead of
displacing anything. The reserved space is identical in every build, so:

\begin{center}
  \emph{the solutions or instructor copy is a pixel-perfect overlay\\
  of the student copy -- same layout, same page breaks, same page count.}
\end{center}

That is why the demo's \file{test.tex} can ship three view PDFs
(section~\ref{sec:assessments}) whose page numbers always agree: page 9 of
the answer key \emph{is} page 9 of the student paper, with answers painted
into its blank space. A quick check after any layout change: a source's
views must all have the same page count. If they differ, something has broken the
contract.

Seeing it live -- this example flips the manual back to overlay mode
locally. The solution box adds no height of its own; it simply paints into
the four workspace lines that follow, exactly where a student would write:

\begin{manexample}
\preserveworkspace
\begin{problem}
  Differentiate $h(x) = x^{5}$.
\end{problem}
\begin{solution}[notitle]
  $h'(x) = 5x^{4}$.
\end{solution}
\workspace[4]
\end{manexample}

In the student build the same four lines are simply blank. Two consequences
follow, both intentional:

\begin{itemize}
  \item \textbf{A solution taller than its workspace overprints the content
    below it.} That collision is the alarm -- it says \emph{leave more room
    for this answer} -- not a bug. Fix it by growing the workspace, never by
    shrinking the answer.
  \item \textbf{Hints are not overlays.} A hint is student-facing help that
    should push content down normally in its build -- it participates in
    layout, unlike the answer key painted over blank space.
\end{itemize}

\subsection{Space the key reclaims: \env{studentonly}}

Some of the room an answer needs is already on the page, spent on something
only the student uses -- most often an empty grid to sketch on, which the key
reprints blank and then draws the real graph underneath. \env{studentonly}
hands that room back:

\begin{manexsource}
\begin{problem}
  Sketch the graph of $f$.
  \begin{solution}          % FIRST: it paints downward
    ...prose, and the graph drawn on a grid...
  \end{solution}
  \begin{studentonly}       % its space is kept, its ink dropped
    \PWgrid{}
  \end{studentonly}
\end{problem}
\workspace[10]
\end{manexsource}

The student copy prints the body; the key prints an empty box of exactly the
same height, so the contract above is untouched -- same page breaks, same page
count -- and the solution now has that height \emph{plus} the workspace to
paint into. Adopting it therefore never makes a previously adequate
\cs{workspace} too small. Under \cs{collapseworkspace} the block vanishes
entirely, like every other reservation.

\textbf{The solution must come first.} A solution written \emph{after} the
block starts below the reclaimed space and paints into the workspace instead,
leaving a blank gap where the grid was. Nothing can detect which you meant, so
nothing warns: the key just looks wrong. Note also that the body is measured
by typesetting it, and in the key that box is then discarded -- so keep
\cs{label}s out of a \env{studentonly} body.

\subsection{Solid mode: the compact answer key}

Sometimes the spaced layout is wrong: an answer-key handout to post online,
or lecture notes (this manual runs in solid mode too). One preamble switch:

\begin{manexsource}
\collapseworkspace  % boxes at natural height; workspace and
                    % \FreshPage collapse to nothing
\preserveworkspace  % back to the default overlay mode
\end{manexsource}

Under \cs{collapseworkspace} the boxes carry their own height and \emph{all} the
reserved room -- every \cs{workspace}, every \cs{FreshPage} break and blank
page -- collapses to nothing: the same source becomes a compact document
with answers inline. The collapse happens only under this explicit,
document-wide switch, never silently per-build; within one mode, every view
of a document lays out identically.

A document may keep \cs{preserveworkspace} written out in its preamble even
though it is the default -- a visible declaration of which mode the document
is in, and a smoke test that the switch pair round-trips.
